\chapter{有计划的商品经济时期1978--1992年}
\label{chap:1978}

\improve[inline]{在本章或下章加入 向社会主义市场经济过渡}

据世界银行数据\footnote{\url{https://data.worldbank.org.cn}}
，1978年的中国人均GDP（现价美元）为156美元，在世界银行所统计的全球130多个国家和地区中，
排名倒数第四，仅高于尼泊尔、几内亚比绍共和国、布隆迪三国，GDP总值149亿美元。是名
副其实的“国穷民弱”。

\section{整体意识形态转变}

\begin{enumerate}
\item 1978年4月5日，中共中央批转公安部《关于全部摘掉“右派”分子帽子的请示报告》，“到1981年，在全面复查的基础上，对错划右派的3改正和落实政策工作全部完成。全国共改正错划右派54万人，占原划右派总数的98\%”\url{http://www.zgdsw.org.cn/n/2013/0125/c244520-20323571.html}。

\item 1978年5月11日，《光明日报》发表胡福明原作后经多次修订的《实践是检验真理的
  唯一标准》，新华社转发。\footnote{\url{http://www.people.com.cn/GB/historic/0510/6267.html}} 次日《人民日报》和《解放军报》同时转载，拉开了全国范围内极左路线失势的序幕，华国锋“两个凡是”理论倒台。
  笔者认为需要说明的是，不考虑本文背后的社会和政治意义，其中的理论本身是正确的，也是马克思一直坚持的，他所关注的实践中的人和社会，正是这种实践的要求，要求理论与时俱进，要求马克思主义理论对其自身的批判和发展。也正因此，马克思抛弃了形而上学的终极实存的哲学体系，也同时抛弃了纯粹的、试图超越人类历史和实践的自然辩证法，他所提倡的是成为运动的、实践的、历史的、辩证的、人的历史唯物主义，并试图改造世界以此让世界更加美好。笔者再次重申，马克思历史唯物主义事实上批判了科学社会主义，详情请见\cref{sec:marxkexue}一节。

\item  1978年11月10日--12月15日，在北京召开中央工作会议，原定议题讨论经济。陈云在东北组率先提出系统地解决历史遗留问题，平凡一系列重大冤假错案。（“到1982年底，约有300多万名干部得到平反。”）邓小平提出“再不实行改革，我们的现代化事业和社会主义事业就会被葬送”\footnote{\url{http://cpc.people.com.cn/GB/64184/64190/65724/4444934.html}}，胡耀邦“认为有些农村体制如“政社合一”就应该改变。” 1978年12月13日，邓小平总结发言《解放思想，实事求是，团结一致向前看》，提到反官僚主义、加强责任制、“允许一部分地区、一部分企业、一部分工人农民，由于辛勤努力成绩大而收入先多一些，生活先好起来”，同时对落后地区和人民给以有力支持等等。

\item 1978年12月18日至22日，中共十一届三中全会在北京召开，继承了之前中央工作会议精神，
  在全中国及历史层面上成为极其重要的一个转折点，另外“邓小平实际上已成为中央领导
  集体的核心。”。\footnote{\url{http://cpc.people.com.cn/GB/64184/64190/65724/4444934.html}}成
  为纠正极左路线、发展市场经济的一个重要转折点，全党工作重点转移到社会主义现代化
  建设上来。

\item 1981年6月27日至29日，中共十一届六中全会在北京召开，会议一致通过《关于建国以来
  党的若干历史问题的决议》，由党和政府系统梳理了这段历史，论述功过，极左势力下台。


\item 1981年9月30日，叶剑英继承和发展了毛泽东、周恩来所提出的“一纲四目”，进一步
  阐明关于台湾回归祖国，实现和平统一的9条方针政
  策。\footnote{\url{http://cpc.people.com.cn/GB/64162/64165/70293/70323/4877134.html}}
  提出“(3)国家实现统一后，台湾可作为特别行政区，享有高度的自治权，并可保留军队。
  中央政府不干预台湾地方事务。(4)台湾现行社会、经济制度不变，生活方式不变，同外国
  的经济、文化关系不变。”1982年，邓小平就叶剑英谈话指出，这实际上就是“一个国
  家，两种制度”。“一国两制”随后也成为香港、澳门回归中国的重要指导思想。邓小平
  在港澳回归问题上立场坚定，行动果决。

\item 1982年9月1日，邓小平在中共第十二次全国人民代表大会开幕式上致开幕词，首次
  提出“有中国特色的社会主义”，开创了全面转向社会主义现代化建设的新局面。
  \begin{quotation}
    我们的现代化建设，必须从中国的实际出发。无论是革命还是建设，都要注意学习和借
    鉴外国经验。但是，照抄照搬别国经验、别国模式，从来不能得到成功。这方面我们有
    过不少教训。把马克思主义的普遍真理同我国的具体实际结合起来，走自己的道路，建
    设有中国特色的社会主义，这就是我们总结长期历史经验得出的基本结论。

    八十年代是我们党和国家历史发展上的重要年代。加紧社会主义现代化建设，争取实现
    包括台湾在内的祖国统一，反对霸权主义、维护世界和平，是我国人民在八十年代的三
    大任务。这三大任务中，核心是经济建设，它是解决国际国内问题的基础。\footnote{\url{http://cpc.people.com.cn/GB/64162/64168/64565/65448/4429495.html}}
  \end{quotation}

\item 1984年10月20日，中共第十二届委员会第三次会议通过《中共中央关于经济体制改革
  的决定》。《决定》中指出“改革计划体制，首先要突破把计划经济同商品经济对立起来
  的传统观念，明确认识社会主义计划经济必须自觉依据和运用价值规律，是\textbf{在公
    有制基础上的有计划的商品经济}。……就总体说，我国实行的是计划经济，即有计划的
  商品经济”。

\item 1986年9月28日，中共十二届六中全会通过《中共中央关于社会主义精神文明建设
  指导方针的决议》。《决议》中指出：
  \begin{quotation}
    我国社会主义现代化建设的总体布局是：以经济建设为中心，坚定不移地进行经济体制
    改革，坚定不移地进行政治体制改革，坚定不移地加强精神文明建设，并且使这几个方
    面互相配合，互相促进。\footnote{\url{http://cpc.people.com.cn/GB/64162/64168/64565/65381/4429515.html}}
  \end{quotation}

\item 1987年10月25日，赵紫阳在中共第十三次代表大会上作了本次大会报告，提出“社会主义
  初级阶段”。
  \begin{quotation}
    我国正处在社会主义的初级阶段。这个论断，包括两层含义。第一，我国社会已经是社
    会主义社会。我们必须坚持而不能离开社会主义。第二，我国的社会主义社会还处在初
    级阶段。我们必须从这个实际出发，而不能超越这个阶段。在近代中国的具体历史条件
    下，不承认中国人民可以不经过资本主义充分发展阶段而走上社会主义道路，是革命发
    展问题上的机械论，是右倾错误的重要认识根源；以为不经过生产力的巨大发展就可以
    越过社会主义初级阶段，是革命发展问题上的空想论，是“左”倾错误的重要认识根
    源。
  \end{quotation}

  在马克思所作的科学共产主义重要纲领性文献《哥达纲领批判》中，将共产主义分为共产
  主义初级阶段（即社会主义）和共产主义高级阶段。“社会主义初级阶段”其实不是“共产
  主义初级阶段”，而是向“共产主义初级阶段”过渡前的带有更多资本产权性质的“前社
  会主义阶段”。

  笔者认为，为方便理解可将此替换为国家资本主义阶段。事实上，不管苏俄或者中国之前
  左的程度如何或它们如何宣称，始终都未曾脱离国家资本主义主体的范畴。在社会主义国
  家阵营的实践中，马恩历史唯物主义中所说的“物质生活条件、生产关系和交换关系的发
  展程度”成为无法跨越的卡夫卡峡谷。这次大会关于生产关系等的认识可说是准确。

  但是，\improve[inline]{item 赵紫阳回忆录中，对一个中心，两个基本点的阐述要补充进来。}

\end{enumerate}

\section{农村改革}

\subsection{包干到户、村社分设及合同订购}

1977年11月15日，安徽省在时任省委第一书记万里的推动下下达了《关于当前农村经
济政策几个问题的规定》（简称“省委六条”），有“尊重生产队的自主权”、“允许和
鼓励社员经营正当的家庭副业，对收回的“自留地”，要按照政策规定如数退还给社
员”等内容，安徽省再兴农业改革。1978年安徽遭遇大旱，省委决定允许借地给农民
种“保命麦”，9月，肥西县山南公社黄花大队与山南公社馆西大队小井庄生产队借此政策
搞起实质“包产到户”，12月凤阳县梨园公社严岗大队小岗生产队决定“包干到
户”。

需要注意的是，虽有省委推动，但不管在党内还是学界中这块内容仍然被认为是农民先导
和主导，学者高王凌将农民主导的软性觉悟称为“反行为”。笔者认为，这种“反行为”实则是中
国各行各业至今仍然具有的普遍现象，有时民间俗称“潜规则”。关于这方面的研究可能会在
相当多的层面上有利于国计民生，希望有相关人士从事和深入此项工作。

% 关于凤阳县“包干到户”，时至今日，仍有种思想认为它错误的将集体性质大幅削弱，笔者
% 不认同这种意见。这种思想没有考虑当时中国实际情况，理论建设在一个空中楼阁之上。正
% 如列宁在二十世纪初期就渴望资本强大、生产力发达的“美国式道路”，希望借此来达到俄
% 国的资本发达，以此作为俄国向社会主义转变的物质基础，后来根据国情只能放弃这一提法，
% 虽经后来斯大林、赫鲁晓夫等努力至苏联解体时苏联仍未真正实现“美国式道路”一样，理
% 论宏观上是能自洽的，实践中却有太多微观细节和问题要面对。

“统购统销”政策需要国家有实力去吸纳全部农村产出以支持“工农业剪刀差”的正常运作：
国家以相对低价格收购农村农副产品，补贴城镇职工，使城镇职工可以保有相对低的工资。
（这种做法在马克思主义理论中其实是提高相对剩余价值。）同时，国家也可在自身维持的
农副产品价格稳定状态中获取相对于自然经济价格的利差（原始积累）去发展工业。毛时代
尤其侧重重工业，邓时代减弱了重工业的扶持力度。但刚经历完多年动荡的我国已经没有力
气再去做全国农副产品的统购统销了。实际上根据薄一波所说，统购统销造成的财政缺口是
巨大的：
\begin{quotation}
  （统购统销）第二年（即1954至1955年度），即赔了2.5506亿元。随着粮食经营费用的增加和购销价格“倒挂”现象的出现，亏损越来越大……1987年达276亿元，1988年突破300亿元，成为国家财政的一大包袱。\pagescite[][281]{boyibo}
\end{quotation}

1983年中央一号文件《当前农村经济政策的若干问题》，从理论、舆论、政策上确立和巩
固了“实行生产责任制，特别是联产承包制；实行政社分设。”原来三级所有的公社、生
产大队和生产队集体组织被取缔，建立乡政府和基层群众性自治的村民委员
会。
\footnote{\url{http://cpc.people.com.cn/GB/64162/64172/85037/85039/6619026.html}}

即使是“包干到户”这种落后小农经济的生产力、生产关系，国家在1984年这个“特大丰收
年”已经无法去完成大批量的“统购统销”了。1984年10月20日，中共十二届三中全会通过
的《中共中央关于经济体制改革的决定》明确转向以城市为重点的整个经济体制改
革。1985年，国家改行“合同订购”制度并在年底要求对订购数量逐年降低，并且“合同订
购”给的农产品价格更低，统购统销的精神纲领仍在。\cite{diyibugaige}国家陆续放开粮
油等农产品价格与经营，实行市场自由购销。

\improve{全国通货膨胀和粮价膨胀}

\begin{quotation}根据农业经济专家\footnote{汪晖此处所指专家为曾任国务院农村发展研究中心发展研究所市场研究室主任等职务的卢迈}的研究，1978~1985年城乡收入的差距是缩小的，从1985年起扩大。1989年到1991年农民收入增长基本停滞，城乡收入差距又恢复到1978年以前的情况。1993年以后，由于国家提高粮食价格、乡镇企业增长快、外出务工人口收入增长等原因，农村收入增长较快，但在城市劳动力大量剩余的情况下，这一势头正在改变。\cite{wangxiandai}
\end{quotation}

三农问题是一个实践上的大难题，涉及工农业剪刀差、剩余劳动力、城乡流动、阶级固化等
问题，它在数年内长期困扰着中国。1992年底，统购统销正式退出历史舞台。时至今日，我
们已经开始了农业工业化的道路，有的旧问题转换了表现形式本质未变，有的旧问题被解决的同
时带来了新问题，我们放在文后再讨论这两个节点吧。\improve[inline]{补充农业工业化
  章节，并替换“文后”为具体章节。}

\subsection{乡镇企业}

1984年3月1日，中共中央、国务院转发农牧渔业部和部党组《关于开创社队企业新局面的报
告》，并同意报告提出的将社队企业名称改为乡镇企业的建议，提出发展多种经营，乡镇企
业是多种经营的重要组成部分。乡镇企业有利于以工补农，吸收农村剩余劳动力，减轻城市
压力等，独立核算、自负盈亏。


\section{城市改革}

1979年4月，中共中央工作会议上广东省委、省政府主要负责人习仲勋、杨尚昆提议发挥广东
优势，“在对外开放商做点文章”，邓小平就此提出“特区”概念。同年7月，中共中央和国
务院正式批准广东福建两省在深圳、珠海、汕头、厦门四个市划出部分地区试办“出口特
区”。1980年5月，中共中央和国务院正式将“出口特区”定名为“经济特区”。人大常委会
此后正式批准此一建议。此后数年持续扩大和增强经济特区的权利和地位。1982年12月，国
务院发布《国务院关于批转<当前试办经济特区工作中若干问题的纪要>的通知》。1984年邓
小平第一次南巡。1月巡视深圳、珠海，2月巡视厦门。同年5月4日，国务院发出《沿海部分
城市座谈会纪要》的通知，确定进一步开放14个沿海港口城市——大连、秦皇岛、天津、烟台、
青岛、连云港、南通、上海、宁波、温州、福州、广州、湛江、北海，扩大厦门经济特区范
围到全岛，进一步开发海南等。随后国务院陆续批准在沿海开放城市及部分地区兴办经济开
发区，引进外资、技术，以开发工业生产新技术和科研为主，减免税收，鼓励出口创
汇。1986年2月发布《国务院关于批转经济特区工作会议纪要的通知》。意识形态上明确
了“在经济上坚持以社会主义经济为领导，允许多种经济成分存在”；在具体举措上则是引
进、利用外资，面向市场经济的国际市场，“以首先搞好基础设施为重点展开基本建设”、
减税或免税等有利于发展市场经济的举措。1988年4月，人大会决定设立海南省经济特区。

1984年10月20日，中共第十二届三中全会通过的《中共中央关于经济体制改革的决定》，明
确了以城市为重点的整个经济体制改革。
\begin{quotation}
  实行政企职责分开以后，要充分发挥城市的中心作用，逐步形成以城市特别是大、中城市
  为依托的，不同规模的，开放式、网络型的经济区……城市政府应该集中力量做好城市
  的规划、建设和管理……城市政府还应该根据国民经济发展的总体要求和当地的条件，做
  好中长期的经济和社会发展规划……使承包责任制在城市生根、开花、结果。
\end{quotation}

\improve[inline]{基建起因、发展、结果}

\section{领导制度改革}
\begin{enumerate}

\item 1980年8月18日，邓小平在中共中央政治局扩大会议上作了《党和国家领导制度的改
  革》的讲话。“陈云同志提出，我们选干部，要注意德才兼备。所谓德，最主要的，就是
  坚持社会主义道路和党的领导。在这个前提下，干部队伍要年轻化、知识化、专业化，并
  且要把对于这种干部的提拔使用制度化。这些意见讲得好。”，要求废除干部领导职务终
  身制，建立退休制度和发扬“社会主义民主制度和党的民主集中制”。
  \url{http://news.12371.cn/2017/03/09/ARTI1489054449475294.shtml}
  \begin{quotation}
    本次会议对中国政治体制中的“官僚主义”、“权力过分集中”的弊端进行了严厉批判，
    将部分过往的政治内容划为封建专制主义，首次提出“改革党和国家领导制度及其他制
    度”的问题。这篇讲话，后来被中共十三大认为是“中国政治体制改革的纲领性文献”，
    也被党内外的主流研究者们奉为研究邓小平政治体制改革思想的经典。\footnote{\url{http://www.yhcqw.com/33/9626.html}}
  \end{quotation}

\item 1982年1月13日，邓小平在中共中央政治局会议上作了《精简机构是一场革命》的重要讲
  话，贯彻和加强了领导制度改革。此后各级机关进行了大规模精简、建立干部离退休制度
  等。1985年到1987年，中国人民解放军也裁军百万。

\item 1982年12月4日全国人大五届五次会议通过并施行的《中华人民共和国宪法》规定全国人
  大委员长、副委员长，中国人民共和国主席、副主席，国务院总理、副总理、国务委员，
  最高人民法院院长，最高人民检察院院长连续连任不得超过两届。彭真在本次大会所作宪
  法修改草案报告中就此项规定说，“这就取消了实际上存在的领导职务的终身制。”


\item 1982年9月召开的中共十二大会议上通过了1982年版的《中国共产党章
  程》\footnote{\url{http://guoqing.china.com.cn/2012-09/05/content_26433894.htm}}，决定设立中央和地方各级顾问委员会。有相当资历和成绩的老党员、老人进入顾问委员会；中顾委接受中央委员会的领导，但具有建议、协调、监督、宣传等权利。

  笔者认为设立中顾委的本意明确，确实如邓小平在9月13日举行的中顾委第一次会议上所
  说
  \begin{quotation}
    中央顾问委员会是个新东西，是根据中国共产党的实际情况建立的，是解决党的中央领
    导机构新老交替的一种组织形式。目的是使中央委员会年轻化，同时让一些老同志在退
    出第一线之后继续发挥一定的作用。……是一种过渡性质的组织形式。
    \footnote{\url{http://www.qstheory.cn/books/2016-08/31/c_1119485398_2.htm}}
  \end{quotation}

  中顾委委员和地方顾问委员会委员人选富有流动性，老顾问委员批量退休，较老的一部
  分一线官员退出一线工作增补进来，并且这种流动性与日俱增。1986年10月，彭真、邓颖超、徐
  向前、聂荣臻四老“全退”，不再担任任何职务。邓小平、陈云、李先念三老“半退”，
  仍担任一定职务。1989年9月，邓小平明确提出十四大以后不再设立顾问委员
  会。1992年10月18日，中共第十四次全国代表大会通过了关于中央顾问委员会工作报告的
  决议，大会同意不再设立中央顾问委员会。中顾委完成了它的历史使命。

\item 就个人的政治利益而言，在邓小平1978年事实上成为核心领导人之际仍有一些与他几
  乎同等资历却未必意见一致的老人。邓小平通过扶持对经济开放持乐观态度的年轻领导班
  子，以此制衡拥有巨大权力的幕后老人——尤其是希望放缓改革开放步伐的老人。通过中顾
  委改革的进行，政坛老人影响力和权力逐步减弱，邓小平成为实际上的“一个婆婆”，为
  其执政思想和方针的实施打下坚实基础。

  老人政治具有相当弊端，进步空间有限，民主成分低，权力集中明显，私心重，人治为上。
  邓小平带着壮士断腕的魄力试图去革除这一弊端，但是在现实的实施上，最后仍然走
  向“一个婆婆”，由统治者个人决定了时代走向。

  就“民主集中制”方面，一个良好的中国“民主集中制”因中国历史和现实问题必然不是
  西方唬烂的那一套民主，但同样因此国情使我们这条道路相当艰辛、带有迷雾、需要更多
  的时间和人去探索。如何在保障集中制为主的基础上加强民主和法治成分，邓小平并未走
  的太远，但他确实在困境下前行了一大步，他大力推动的党和领导制度改革的功绩不能抹
  杀，不能对其过于苛责。
\end{enumerate}

\section{财政体制改革}

\subsection{财政体制改革的深层逻辑}

2008年3月18日，时任国务院总理温家宝在十一届全国人大一次会议的记者见面会上说：“其
实一个国家的财政史是惊心动魄的。如果你读它，会从中看到不仅是经济的发展，而且是社
会的结构和公平正义。”

\begin{quotation}
  回顾40年改革长路，是什么在推动中国财政改革？其深层逻辑又是什么？我们发现，从过
  去到现在，这一问题根本上可以归结为一点，那就是公共风险的变化。\textbf{财政改革
    往往不是在先见之明的制度设计基础上进行的，而是在公共风险暴露与加剧时，不得不
    做出的选择。}

  　防范和化解公共风险是财政改革的原动力。从历史上看，制度变迁无一不是公共风险与
  危机推动的结果，而制度变迁的突破口常常是在财政。改革开放以来，我国财政改革实质
  上都是遵循公共风险变化的逻辑而推进的，其变化的脉络是从“\textbf{家贫国穷}”的风
  险到“\textbf{机会不均}”的风险，再到\textbf{全球公共风险}，这也是我国主要公共
  风险的昨天、今天和明天。
\end{quotation}

\subsection{就业制度改革}
1980年8月17日，中共中央发文《中共中央关于转发全国劳动就业会议文件的通知》，
指出“要积极创造条件，在国家统筹规划和指导下，实行劳动部门介绍就业、自愿组织起
来就业和自谋职业相结合”的三结合方针，企业公开招工并可以增减人员。扭转过往片面
强调“大”和“公”的错误以及“铁饭碗”的状况。1983年2月22日，劳动人事部发文《劳
动人事部关于积极试行劳动合同制的通知》，积极实行劳动合同制，打破“大锅
饭”、“铁饭碗”，调动人民积极性，充分调动人们的社会主义积极性，解放生产
力，“（全民所有制单位及集体所有制单位）经过若干步骤，最终达到所有职工都实行劳
动合同制。”
\footnote{\url{http://fgcx.bjcourt.gov.cn:4601/law?fn=chl008s060.txt}}
\begin{quotation}
  国务院于1986年7月发布了关于国营企业劳动制度改革的四项暂行规定……然而,这一改
  革措施只限于国营企业中新招收的工人,并没有根本动摇原有的固定工制度。到1991年
  底,全民所有制单位合同制工人只有1449万人,占全民所有制单位职工总数的14\%。\cite{laodongzhiduyuanfang}
\end{quotation}

\improve[inline]{就业制度改革，以后参看程连升的《中国五十年反失业政策研究（1949--1999）》}

\subsection{扩大企业经营管理自主权}

1978年，四川省委、人民政府在6个地方国营工业企业开展扩大工业企业经营管理自主权试
点，在1979年1月扩大至100家工业企业和40家商业企业。1979年7月后，国务院出台一系列
文件要求全国进行扩大企业自主权试点，1984年5月10日，国务院颁布了《国务院关于进一
步扩大国营工业企业自主权的暂行规定》。1984年10月20日，中共第十二届三中全会通过
《中共中央关于经济体制改革的决定》，在宏观理论层面上确定了城市企业学习农村，大
力发展建立以承包为主的多种形式的经济责任制；增强企业活力，减弱政府对经济的各种
管控，政企职责分开，确立企业主体地位；尊重价值规律及建立合理价格体系等。

\improve[inline]{86年企业转租}

\improve[inline]{1992年，国务院颁布《全民所有制工业企业
  转换经营机制条例》。 应该算到下一章节。}


\subsection{税制改革}

1980年9月10日，第五届全国人民代表大会第三次会议通过了《中华人民共和国中外合资经
营企业所得税法》和《中华人民共和国个人所得税法》。1981年12月13日，第五届全国人民
代表大会第四次会议通过了《中华人民共和国外国企业所得税法》。中国涉外税制初步建立。

1980年2月，国务院颁发了《关于实行"划分收支，分级包干"的财政管理体制的规定》，全面
推行包干式财政管理体制，扩大地方财政权限，俗称“分灶吃饭”，改变了过去“统收统
支”的财政体制。这一政策实行至1984年。

1983年进行了国有企业上交利润转上交税收的第一步改革，小型国有企业自负盈亏。1984年
进行了第二步利改税改革，进一步分解增设多种税种，由税利并存逐步过渡到完全的以税代利。

1985年3月21日，国务院发布《关于实行“划分税种、核定收支、分级包干”的财政管理体制
的规定》，进一步改革包干式财政管理体制，“一定五年不变”。

税制改革从此成为中央与地方政府的博弈场。依马骏所说，原“分灶吃饭”税制
中，“20\%共享收入归地方的一刀切计算办法使富省多有结余，穷省多有赤字”，因此促成
了这次改革。新税制改革的弊端是：
\begin{quotation}
  自1978年底以来，政府财政收入占国民生产总值的比例及中央财政收入占全部财政收入的
  比重快速地下降。

  （地方政府）利用各种减免税来调节其实际税收努力……地方政府对有效税率及税基的控
  制还表现在其对财政收入渠道的控制……地方政府预算内收入比重的下降和预算外收入上
  升，中央政府控制全国财政收入与支出的能力相应下降。

  由于地方政府在事实上控制了有效税率和税基，中央政府越来越多地依赖于一些不规范的
  政策工具来控制地方向中央上缴的财政收入。\cite{majuncaigai}
  \improve[inline]{分税制也要借鉴这篇文章}

\end{quotation}

1985年，所有国家基本建设投资由财政无偿拨款改为向中国人民建设银行贷款，并要还付本
息，即“拨改贷”。信托投资公司兴起。
\begin{quotation}
  1984年国家出台拨款改贷款政策后，意味着1984年以后新成立的国企，一出生就是
100\%的负债率。\cite{bogaidaizhaizhuangu}

国有企业高负债率的直接原因是国家在1983年实行了拨改贷，对国有企业的投资由财政拨款
改为向银行的贷款。国有企业高负债率的间接原因则是国有企业的预算软约束。当国有企业
发生亏损或经营困难时，国家由传统的财政拨款直接支持改为间接地由银行的低息贷款来支
持，国有企业贷了款后经营状况没有改善，借的款越来越多，负债的比率也就越来越高。\cite{linyifuzhai}
\end{quotation}

\improve[inline]{由“拨改贷”到“债转
  股”\url{http://www.hprc.org.cn/gsyj/jjs/jjyxs/201702/t20170220_388850.html}}

 财政部官网转载刘邦驰、马韵《三十年税制改革的回顾与发展趋向》\footnote{\url{http://www.mof.gov.cn/zhuantihuigu/czgg0000_1/czggcz/200811/t20081106_88385.html}}
  ，将1978年--1993年这一时期的税制改革称为“有计划商品经济时期的税制改革”，涉外
  税制、利改税(“新中国成立后实行了 30 多年的国营企业向国家上缴利润的制度改为缴纳
  企业所得税和调节税”)、个人所得税等税务改革。
  \begin{quotation}
    这一时期全面改革了工商税制，建立了涉外税制，彻底摒弃了“非税论”和“税收无
    用论”的观点，恢复和开征了一些新税种，从而使我国税制逐步转化为多税种、多环
    节、多层次的复合税制，税收调节经济的杠杆作用日益加强。\footnote{\url{http://finance.ce.cn/sub/2013zt/zgszgg/hgyzw/201304/18/t20130418_199980.shtml}}
  \end{quotation}


陈云 \url{http://cpc.people.com.cn/GB/69112/83035/83317/83597/5738412.html}《经济形势与经验教训 （一九八〇年十二月十六日）》

  1981年1月7日，国务院《关于加强市场管理、打击投机倒把、走私贩私活动的指示》
  1982年4月10日政治局会议讨论《中共中央、国务院关于打击经济领域中严重犯罪活动的
  决定》。4月13日，中共中央、国务院发布了《关于打击经济领域中严重犯罪活动的决定》。
  1983年9月15日，人民日报发表评论员文章《必须依法从重从快惩处经济罪犯》。1983年
  月12日，中共第十二次全国代表大会发出《中共中央关于整党的决定》。

  1983年8月，全国政法工作会议公布《中共中央关于严厉打击刑事犯罪活动的决定》，全
  国从重从速严打刑事犯罪。



  从1978--1992年是中国左右思想交错争锋的一段时间。\improve{未完}



\improve[inline]{中央与地方张力，税务历史、}
% \begin{verbatim}
% https://www.modernchinastudies.org/us/issues/past-issues/45-mcs-1995-issue-1/253-2011-12-29-11-30-06.html
% 中央与地方财政关系的改革

% 南方周末：重议分税制——分钱还是分权？ https://www.guancha.cn/economy/2013_08_15_165993.shtml

%  分税制的是与非 http://liushangxi.blog.caixin.com/archives/60712

% http://www.sss.net.cn/112/65928.aspx
% 刘尚希:财政改革四十年的深层逻辑

% 回顾40年改革长路，是什么在推动中国财政改革？其深层逻辑又是什么？我们发现，从过去到现在，这一问题根本上可以归结为一点，那就是公共风险的变化。财政改革往往不是在先见之明的制度设计基础上进行的，而是在公共风险暴露与加剧时，不得不做出的选择。

% 　防范和化解公共风险是财政改革的原动力。从历史上看，制度变迁无一不是公共风险与危机推动的结果，而制度变迁的突破口常常是在财政。改革开放以来，我国财政改革实质上都是遵循公共风险变化的逻辑而推进的，其变化的脉络是从“家贫国穷”的风险到“机会不均”的风险，再到全球公共风险，这也是我国主要公共风险的昨天、今天和明天。

% http://www.cssn.cn/zx/bwyc/201806/t20180626_4421303.shtml
% 中央、地方与企业：财政体制改革40年

% http://theory.people.com.cn/GB/49154/49155/8094994.html

% http://www.pkulaw.cn/fax/czf/contents/class02_3_1r.htm

% http://www.hprc.org.cn/gsyj/jjs/jjyxs/201702/P020170220376223183282.pdf
% 贷改投

% http://www.ccswf.org.tw/files/7100/18/%E5%AD%90%E9%A1%8C3-2%E8%94%A1%E7%A6%BE.pdf

% 建国以来最高通货膨胀是什么时候？
% \end{verbatim}

% ( 4) 共产主义社会高级阶段。 中国共产党第十三次全国代表大会的报告中指出, 社会主义初级阶% 段 “ 不是泛指任何国家进入社会主义社会都会经历的起始阶段, 而是特指我国在生产力落后、 商品% 经济不发达条件下建设社会主义必然要经历的特定阶段 。” [6 ] % 发达资本主义国家在无产阶级夺取 出自《正确认识和对待“共产主义渺茫论”》
\improve[inline]{\url{http://theory.people.com.cn/GB/148980/16813293.html}，邓小
  平南巡的三个坚定不移。}


%%% Local Variables:
%%% mode: latex
%%% TeX-master: "../main"
%%% End:
